\section{Definición y clasificación de los Agents: percepción + acción}

La introducción mostró que los cuatro materiales apuntan a los Agents; esta sección primero asienta la base conceptual y responde una pregunta concreta: qué sistema cuenta como Agent y cuál es solo un chatbot. La descripción del video da una definición concisa: un Agent es un sistema inteligente capaz de interactuar con el entorno, con capacidad de percepción y capacidad de acción. La percepción incluye observar el entorno, leer la retroalimentación e interpretar el contexto; la acción incluye llamar herramientas, generar salidas, controlar interfaces y modificar variables. El ciclo mínimo es «observar → decidir → actuar».

\begin{figure}[H]
\centering
\includegraphics[width=0.96\textwidth]{agent-definition-loop.png}
\caption{Definición mínima de un Agent: percibir el entorno, decidir, actuar y volver a leer la retroalimentación. Diagrama conceptual propio, elaborado a partir de la conversación de 00:02:00--00:14:50.}
\end{figure}

\begin{knowledgebox}{Lectura de la figura: diferencia entre un Agent y un Chatbot}
Un Chatbot principalmente responde a la entrada; un Agent tiene que cambiar el estado del entorno. En cuanto el sistema empieza a llamar herramientas, leer la retroalimentación y ajustar su plan, deja de ser un sistema conversacional y pasa a ser un sistema de agente.
\end{knowledgebox}

\subsection{Cuatro tipos de Agent}

Con la definición mínima establecida, esta subsección los clasifica según el entorno y la forma de producto. El programa divide los Agents en Coding Agent, Search Agent, Tool-use Agent y Computer-use Agent. Esta clasificación no es una clasificación académica rigurosa, sino una clasificación por producto y entorno: código, búsqueda, herramientas e interfaz de computadora corresponden a distintas affordance y formas de verificación.

\begin{figure}[H]
\centering
\includegraphics[width=0.96\textwidth]{agent-types-map.png}
\caption{Mapa de tipos de Agent: Coding, Search, Tool-use y Computer-use corresponden a entornos distintos. Diagrama conceptual propio, elaborado a partir de la conversación de 00:02:00--00:14:50.}
\end{figure}

\begin{knowledgebox}{Términos clave: tipos de Agent}
\begin{tabularx}{\textwidth}{p{0.22\textwidth}p{0.34\textwidth}X}
\toprule
Tipo & Tareas representativas & Dificultades \\
\midrule
Coding Agent & Autocompletado de código, refactorización, depuración, modificaciones de PR & Contexto del repositorio de código, pruebas, modificaciones de largo alcance. \\
Search Agent & Recuperación, consolidación, informes de investigación & Fiabilidad de las fuentes, citas, eliminación de duplicados y synthesis. \\
Tool-use Agent & Llamar API, funciones, bases de datos y herramientas de negocio & Selección de herramientas, generación de parámetros, recuperación ante errores. \\
Computer-use Agent & Operar el navegador, la GUI y flujos entre aplicaciones & Grounding visual, seguimiento del estado, seguridad de permisos. \\
\bottomrule
\end{tabularx}
\end{knowledgebox}

\subsection{Resumen de la sección}

La definición de un Agent puede ser muy sencilla, pero llevarla a la práctica es muy complejo. La diferencia esencial entre los distintos tipos de Agent está en que enfrentan entornos, herramientas y retroalimentación distintos.
